\section{Introduction: the eight questions behind a yes vote}
\label{sec:intro}

A bill does not pass because it is correct. It passes because enough members of each
chamber can answer, to their own satisfaction and to a constituent's, the small set
of questions that stand between a draft and a statute. This paper sets out those
questions for H. R. 9510, the Verification Before Generation in Physical AI Oncology
Trials Act of 2026, and answers each with evidence the member can cite
\cite{Kawchak2026HR9510Bill}.

The subject is unfamiliar to most members, so the terms are defined before they are
used. Physical AI is the integration of embodied robots with advanced artificial
intelligence for direct patient care in a clinical trial
\cite{Kawchak2026NationalPlatformPhysicalAI}. Robot-patient interaction code is the
machine-generated software that directs a robot's behavior at the moment it touches a
patient. Verification, validation, and uncertainty quantification, abbreviated VVUQ,
is an examination that asks whether the action was built correctly, whether it is the
right action for this patient, and how confident the system is. Verification before
generation is the bill's central rule: the examination must clear before the code is
written or run, not after. The examination is a suite of ten gates, each bound to a
published safety standard, and each action it reviews is resolved to one of three
outcomes, ACCEPT and act under audit, ESCALATE to a qualified human, or BLOCK before
the patient. The bill writes this rule into a new section 515D of the Federal Food,
Drug, and Cosmetic Act. \autoref{tab:terms} collects these and the budget terms a
member needs.

\needspace{12\baselineskip}\begin{table}[H]
\footnotesize
\begin{tabularx}{\textwidth}{@{}L{4.6cm} Y@{}}
\toprule
\textbf{Term} & \textbf{Plain-language meaning for a legislator}\\
\midrule
Physical AI & Embodied robots guided by advanced AI for direct patient care in a trial.\\
Robot-patient interaction code & The software that directs a robot at the moment it contacts a patient.\\
VVUQ & Automated check if an action is built right, is for patient, and how certain system is.\\
Verification before generation & The rule that the check must clear before code is generated or executed.\\
Ten gates; ACCEPT, ESCALATE, BLOCK & The ten-part check and its three outcomes; three gates are hard catastrophe predicates that must always hold.\\
Section 515D & The new section of the FD\&C Act the bill adds (21 U.S.C. 360e-5).\\
Authorization versus appropriation & Authorization sets a ceiling Congress may later fund; appropriation is actual money. \\
PAYGO and CutGo & Budget rules requiring that new mandatory spending be offset; this bill creates none.\\
\bottomrule
\end{tabularx}
\caption{The terms a member needs before reading the framework, defined before first
substantive use.}
\label{tab:terms}
\end{table}

The argument is an ordered path. The eight questions run from whether a statute is
needed at all, through authority, safety, and the budget score, to the people the
bill helps, the politics of the two caucuses, the coalition outside the chamber, and
finally the path through both chambers. \autoref{fig:intro-path} traces that path;
each later section answers one question and closes on the provision, score
line, or settling evidence.

\begin{figure}
\begin{psgfig}
\node[psgone] (q1) at (0,0) {1 Mandate};
\node[psgtwo] (q2) at (3.9,0) {2 Authority};
\node[psgact] (q3) at (7.8,0) {3 Safety};
\node[psgthree] (q4) at (11.7,0) {4 Fiscal};
\node[psgtwo] (q5) at (11.7,-2.0) {5 Constituents};
\node[psgevid] (q6) at (7.8,-2.0) {6 Bipartisan};
\node[psgthree] (q7) at (3.9,-2.0) {7 Coalition};
\node[psgact] (q8) at (0,-2.0) {8 Passage};
\draw[psgedge] (q1)--(q2); \draw[psgedge] (q2)--(q3); \draw[psgedge] (q3)--(q4);
\draw[psgedge] (q4)--(q5);
\draw[psgedge] (q5)--(q6); \draw[psgedge] (q6)--(q7); \draw[psgedge] (q7)--(q8);
\end{psgfig}
\vspace{0.1cm}
\psgcaption{\textbf{Figure 1.} The eight passage questions as one ordered path, from
the mandate for a statute to the path through chambers.}
\label{fig:intro-path}
\end{figure}

The objective of the framework is not to win an argument. It is to assemble a
majority in the House and the votes to pass the Senate, by giving each member an
answer they can defend in a markup.

\clearpage
